\documentclass[11pt]{article}

\usepackage[margin=1in]{geometry}
\usepackage{graphicx}
\usepackage{amsmath, amssymb}
\usepackage{braket}
\usepackage[hidelinks]{hyperref}

\title{ECE 523 Quantum Computing Portfolio}
\author{Manny Serrano Campa}
\date{Fall 2026}

\begin{document}

\maketitle
\tableofcontents
\newpage

\section*{Introduction}
\addcontentsline{toc}{section}{Introduction}

This portfolio contains reading notes, exercises and glossary terms related to Quantum Computing, developed throughout the Fall 2026 semester as part of my ECE 523 Quantum Computing course. 

Exercises were taken from the Exercise Set pdf constructed by Prof. Natalie Klco. // insert link here. 


% Add your introduction as the portfolio develops.

\part{Reading Notes}

% Add reading notes here.

\part{Exercises}

% Match the exercise set: I, II, III, ...
\setcounter{section}{0}
\renewcommand{\thesection}{\Roman{section}}
\renewcommand{\thesubsection}{\Alph{subsection}}
\renewcommand{\thesubsubsection}
    {\thesection.\thesubsection.\arabic{subsubsection}}

\section{Warm-Up and Preliminaries}

\subsection{Tensor Products}

\subsubsection{Parameters Needed to Specify a Quantum State}
\label{ex:state-parameters}

We count the independent real parameters needed to specify a
pure physical state. Normalization constrains the amplitudes,
and an overall phase does not distinguish physical states.

\paragraph{Starting with one qubit.}
A single qubit belongs to a two-dimensional Hilbert space and
can be written as
\[
\ket{\psi}=a\ket{0}+b\ket{1},
\qquad a,b\in\mathbb{C}.
\]
Each complex amplitude contains two real numbers, so we begin
with four real parameters. Normalization requires
\[
|a|^2+|b|^2=1,
\]
which removes one independent real parameter because the total
measurement probability must equal one.

An overall phase removes one more parameter. To see why,
multiplying the entire state by $e^{i\gamma}$ leaves the
probability of any projective measurement outcome unchanged:
\[
\left|\bra{u}e^{i\gamma}\ket{\psi}\right|^2
=
\left|e^{i\gamma}\braket{u|\psi}\right|^2
=
\left|\braket{u|\psi}\right|^2.
\]
This freedom applies to the global phase; the relative phase
between $a$ and $b$ can affect interference and must be retained.
Therefore, one qubit requires
\[
4-1-1=2
\]
independent real parameters.

\paragraph{A separable state of $n$ qubits.}
For a fully separable pure state,
\[
\ket{\Psi}
=
\ket{\psi_1}\otimes\ket{\psi_2}\otimes\cdots
\otimes\ket{\psi_n},
\]
we can specify each qubit individually. Since each qubit
requires two real parameters, the total is
\[
\boxed{N_{\mathrm{separable}}=2n.}
\]
There are no additional physical parameters associated with
the individual global phases: multiplying each factor by
$e^{i\gamma_j}$ only multiplies the full product by the
unobservable global phase $e^{i\sum_j\gamma_j}$.

\paragraph{A general state of $n$ qubits.}
The combined Hilbert space has dimension $2^n$. A general
pure state can therefore be expanded as
\[
\ket{\Psi}
=
\sum_{x=0}^{2^n-1}c_x\ket{x},
\]
where $\ket{x}$ labels a computational basis state using
the $n$-bit binary representation of $x$.

There are $2^n$ complex amplitudes, giving $2\cdot 2^n$
real parameters before imposing physical constraints.
The normalization condition
\[
\sum_{x=0}^{2^n-1}|c_x|^2=1
\]
removes one parameter, and the global phase removes another.
Thus,
\[
\boxed{N_{\mathrm{general}}=2^{n+1}-2.}
\]
This is the parameter count for an unrestricted pure state,
including a generic entangled state. 

\paragraph{}The difference in scaling comes from whether the joint
amplitudes factor into individual qubit amplitudes.
A separable state can be described one qubit at a time,
so its parameter count grows linearly with $n$.
A general state requires amplitudes for the full joint basis,
so its parameter count grows exponentially.

\paragraph{$d$-level quantum degrees of freedom}
A single $d$-level system has $d$ complex amplitudes.
After removing normalization and global phase, it requires
$2d-2$ independent real parameters. For $n$ fully separable
$d$-level systems,
\[
\boxed{N_{\mathrm{separable}}=n(2d-2)=2n(d-1).}
\]

For a general joint pure state, the Hilbert space has
dimension $d^n$. Applying the same counting argument gives
\[
\boxed{N_{\mathrm{general}}=2d^n-2.}
\]

As a consistency check, setting $d=2$ recovers the qubit
results. Setting $n=1$ makes both counts equal to $2d-2$,
as expected: with only one subsystem, there is no distinction
between a product state and a general state arising from
entanglement between subsystems.

% Add subsequent exercises here as we work through them.

\subsubsection{Tensor Product of Three Identical Qubit States}
\label{ex:three-qubit-product}

The notation $\ket{\psi}^{\otimes 3}$ means taking the tensor
product of $\ket{\psi}$ with itself three times. For
\[
\ket{\psi}
=
\frac{\ket{0}+\sqrt{2}\ket{1}}{\sqrt{3}},
\]
the joint state is
\[
\ket{\psi}^{\otimes 3}
=
\frac{1}{3\sqrt{3}}
\left(\ket{0}+\sqrt{2}\ket{1}\right)
\otimes
\left(\ket{0}+\sqrt{2}\ket{1}\right)
\otimes
\left(\ket{0}+\sqrt{2}\ket{1}\right).
\]

The tensor product distributes over addition. Starting with
the first two factors gives
\[
\ket{\psi}^{\otimes 2}
=
\frac{1}{3}
\left(
\ket{00}
+\sqrt{2}\ket{01}
+\sqrt{2}\ket{10}
+2\ket{11}
\right).
\]
Taking the tensor product with the third qubit then gives
\[
\begin{aligned}
\ket{\psi}^{\otimes 3}
=
\frac{1}{3\sqrt{3}}\Big(
&\ket{000}
+\sqrt{2}\ket{001}
+\sqrt{2}\ket{010}
+2\ket{011}\\
&+\sqrt{2}\ket{100}
+2\ket{101}
+2\ket{110}
+2\sqrt{2}\ket{111}
\Big).
\end{aligned}
\]
Each coefficient is the product of the three individual
qubit amplitudes. Thus, each occurrence of a $1$ contributes
a factor of $\sqrt{2}$ inside the parentheses. For example,
$\ket{101}$ has two ones, giving a coefficient of
$(\sqrt{2})^2=2$.

Using the computational basis in binary order,
\[
\left(
\ket{000},\ket{001},\ket{010},\ket{011},
\ket{100},\ket{101},\ket{110},\ket{111}
\right),
\]
we obtain the column vector
\[
\boxed{
\ket{\psi}^{\otimes 3}
=
\frac{1}{3\sqrt{3}}
\begin{pmatrix}
1\\
\sqrt{2}\\
\sqrt{2}\\
2\\
\sqrt{2}\\
2\\
2\\
2\sqrt{2}
\end{pmatrix}
}
\]

There are eight entries because the three-qubit Hilbert
space has dimension $2^3=8$. As a normalization check,
the squared amplitudes sum to
\[
\frac{1}{27}
\left(1+2+2+4+2+4+4+8\right)=1.
\]
Although the expansion contains all eight computational
basis states, the state is separable since it was constructed
as a tensor product of three individual qubit states.
A superposition of computational basis states does not,
by itself, imply entanglement.

\subsubsection{Tensor Products and Reordering Subsystems}
\label{ex:tensor-product-algorithm}

The tensor product creates a larger vector space by combining
the vector spaces of the individual systems. We can start
with two vectors:
\[
\ket{a}=
\begin{pmatrix}a_0\\a_1\end{pmatrix},
\qquad
\ket{b}=
\begin{pmatrix}b_0\\b_1\end{pmatrix}.
\]
To take their tensor product, we multiply each entry of
$\ket{a}$ by the entire vector $\ket{b}$:
\[
\ket{a}\otimes\ket{b}
=
\begin{pmatrix}
a_0\ket{b}\\
a_1\ket{b}
\end{pmatrix}
=
\begin{pmatrix}
a_0b_0\\
a_0b_1\\
a_1b_0\\
a_1b_1
\end{pmatrix}.
\]
So two vectors with two entries each give a joint vector
with four entries.

\paragraph{Implementing the tensor product.}
For vectors of lengths $d_1$ and $d_2$, we first create
an output vector with $d_1d_2$ entries. Then we go through
each entry of the first vector and multiply it by each
entry of the second vector, keeping that order.

Using indices starting at zero, this is
\[
c_{i d_2+j}=a_i b_j,
\qquad
0\leq i<d_1,\quad 0\leq j<d_2.
\]
The index $i$ stays fixed while $j$ goes through the second
vector. Then we move to the next value of $i$.

For matrices, we do the same thing: replace each entry
of $A$ with that entry times the entire matrix $B$:
\[
A\otimes B=
\begin{pmatrix}
A_{00}B & \cdots & A_{0,n-1}B\\
\vdots & \ddots & \vdots\\
A_{m-1,0}B & \cdots & A_{m-1,n-1}B
\end{pmatrix}.
\]
If $A$ is $m\times n$ and $B$ is $p\times q$, the result
is an $mp\times nq$ matrix. Its entries are
\[
(A\otimes B)_{ip+k,\;jq+\ell}=A_{ij}B_{k\ell}.
\]

\paragraph{How the basis is ordered.}
For the two-qubit example, the entries correspond to
\[
(\ket{00},\ket{01},\ket{10},\ket{11}).
\]
The rightmost bit changes first. This is the same as
counting in binary from $0$ to $3$.

More generally, the basis state
$\ket{i}_1\otimes\ket{j}_2$ has index
\[
x=i d_2+j.
\]
We go through every state of the second space before
moving to the next state of the first space. If both
spaces use binary ordering, we join their bit strings
to label the joint state.

\paragraph{Building the reordering matrices.}
To build each matrix, we look at where every input
basis state goes. The input determines the column,
and the output determines the row. We put a $1$ at
that position and zeros everywhere else.

If the input index $x$ goes to $f(x)$, then
\[
U=\sum_x\ket{f(x)}\bra{x}.
\]
Below, all row and column indices start at zero.
A listed pair $(r,c)$ means $U_{r,c}=1$, so it also
specifies the outer-product term $\ket{r}\bra{c}$.

\paragraph{Reversing three qubits.}
For the order $(1,2,3)\rightarrow(3,2,1)$, we swap
the first and third bits:
\[
U_{\mathrm{rev}}\ket{abc}=\ket{cba}.
\]
The input index is $4a+2b+c$, and the output index
is $4c+2b+a$. For example,
\[
\ket{001}\rightarrow\ket{100},
\]
so column $1$ has its $1$ in row $4$.

Applying this to all eight basis states gives
\[
U_{\mathrm{rev}}=
\begin{pmatrix}
1&0&0&0&0&0&0&0\\
0&0&0&0&1&0&0&0\\
0&0&1&0&0&0&0&0\\
0&0&0&0&0&0&1&0\\
0&1&0&0&0&0&0&0\\
0&0&0&0&0&1&0&0\\
0&0&0&1&0&0&0&0\\
0&0&0&0&0&0&0&1
\end{pmatrix}.
\]
The eight nonzero entries are
\[
\begin{aligned}
(r,c)\in\{&
(0,0),(4,1),(2,2),(6,3),\\
&(1,4),(5,5),(3,6),(7,7)\}.
\end{aligned}
\]

\paragraph{Reordering four qubits.}
For the order $(1,2,3,4)\rightarrow(2,3,1,4)$,
the basis states change as
\[
U_{\mathrm{cycle}}\ket{abcd}=\ket{bcad}.
\]
The input and output indices are
\[
x_{\mathrm{in}}=8a+4b+2c+d,
\qquad
x_{\mathrm{out}}=8b+4c+2a+d.
\]
For example,
\[
\ket{1010}\rightarrow\ket{0110},
\]
which gives a $1$ in row $6$, column $10$.

Since the fourth bit stays the same, we can write
the $16\times16$ matrix using $2\times2$ blocks:
\[
U_{\mathrm{cycle}}=
\begin{pmatrix}
I&0&0&0&0&0&0&0\\
0&0&0&0&I&0&0&0\\
0&I&0&0&0&0&0&0\\
0&0&0&0&0&I&0&0\\
0&0&I&0&0&0&0&0\\
0&0&0&0&0&0&I&0\\
0&0&0&I&0&0&0&0\\
0&0&0&0&0&0&0&I
\end{pmatrix},
\qquad
I=\mathbb{I}_2,\quad 0=0_{2\times2}.
\]
Each identity block keeps the fourth bit unchanged.
The sixteen nonzero entries are
\[
\begin{aligned}
(r,c)\in\{&
(0,0),(1,1),(4,2),(5,3),\\
&(8,4),(9,5),(12,6),(13,7),\\
&(2,8),(3,9),(6,10),(7,11),\\
&(10,12),(11,13),(14,14),(15,15)\}.
\end{aligned}
\]

\paragraph{Swapping two three-level systems.}
Now each subsystem has three basis states, labeled
$0$, $1$, and $2$. The joint basis is
\[
(\ket{00},\ket{01},\ket{02},
 \ket{10},\ket{11},\ket{12},
 \ket{20},\ket{21},\ket{22}).
\]
We want
\[
U_{\mathrm{swap}}\ket{ab}=\ket{ba}.
\]
The input index is $3a+b$, and the output index is
$3b+a$. This gives
\[
U_{\mathrm{swap}}=
\begin{pmatrix}
1&0&0&0&0&0&0&0&0\\
0&0&0&1&0&0&0&0&0\\
0&0&0&0&0&0&1&0&0\\
0&1&0&0&0&0&0&0&0\\
0&0&0&0&1&0&0&0&0\\
0&0&0&0&0&0&0&1&0\\
0&0&1&0&0&0&0&0&0\\
0&0&0&0&0&1&0&0&0\\
0&0&0&0&0&0&0&0&1
\end{pmatrix}.
\]
The nine nonzero entries are
\[
\begin{aligned}
(r,c)\in\{&
(0,0),(3,1),(6,2),\\
&(1,3),(4,4),(7,5),\\
&(2,6),(5,7),(8,8)\}.
\end{aligned}
\]

\paragraph{Checking the results.}
Each matrix has exactly one $1$ in every row and column.
This means its columns are normalized and orthogonal,
so
\[
U^\dagger U=\mathbb{I}.
\]
Thus, all three transformations are unitary. They
rearrange the amplitudes while preserving the total
probability.

We can also check how many times we must apply each
transformation to get back to the original state.
Reversing three qubits twice or swapping two qutrits
twice restores the original order:
\[
U_{\mathrm{rev}}^2=\mathbb{I}_8,
\qquad
U_{\mathrm{swap}}^2=\mathbb{I}_9.
\]
The four-qubit transformation cycles the first three
positions, so it takes three applications:
\[
U_{\mathrm{cycle}}^3=\mathbb{I}_{16}.
\]

\setcounter{subsection}{2}
\subsection{Hermitian and Unitary Matrices}

\setcounter{subsubsection}{1}
\subsubsection{Time Evolution for a Time-Independent Hamiltonian}

We start with the Schr\"odinger equation,
\[
i\hbar\frac{\partial}{\partial t}\ket{\psi(t)}
=
\hat{H}\ket{\psi(t)},
\]
where $\hat{H}$ does not depend on time. We want to find
the operator that takes an initial state at time $t_0$
to its state at time $t$:
\[
\ket{\psi(t)}=U(t,t_0)\ket{\psi(t_0)}.
\]

Substituting this into the Schr\"odinger equation gives
\[
i\hbar\frac{\partial U(t,t_0)}{\partial t}
\ket{\psi(t_0)}
=
\hat{H}U(t,t_0)\ket{\psi(t_0)}.
\]
The initial state $\ket{\psi(t_0)}$ is fixed, so the
time derivative acts only on $U(t,t_0)$. Since this
equation must hold for any initial state, the operators
must satisfy
\[
\frac{\partial U(t,t_0)}{\partial t}
=
-\frac{i}{\hbar}\hat{H}U(t,t_0).
\]

This has the same form as the differential equation
$\frac{df}{dt}=af$ with constant $a$, whose solution
is an exponential. Here, the constant is the operator
$-i\hat{H}/\hbar$. We also need
\[
U(t_0,t_0)=\mathbb{I},
\]
because evolving through zero time must leave the
state unchanged. Together, these give
\[
\boxed{
U(t,t_0)=
e^{-\frac{i}{\hbar}\hat{H}(t-t_0)}.
}
\]
Therefore, the state at time $t$ is
\[
\boxed{
\ket{\psi(t)}
=
e^{-\frac{i}{\hbar}\hat{H}(t-t_0)}
\ket{\psi(t_0)}.
}
\]

\paragraph{Checking the solution.}
Since $\hat{H}$ is time independent, differentiating
the matrix exponential gives
\[
\frac{\partial U}{\partial t}
=
-\frac{i}{\hbar}\hat{H}
e^{-\frac{i}{\hbar}\hat{H}(t-t_0)}
=
-\frac{i}{\hbar}\hat{H}U,
\]
which is the equation we needed to solve. At $t=t_0$,
the exponent is zero, so $U(t_0,t_0)=\mathbb{I}$.

For a short time interval $\Delta t$, the expansion is
\[
U(t_0+\Delta t,t_0)
=
\mathbb{I}
-\frac{i}{\hbar}\hat{H}\Delta t
+\mathcal{O}(\Delta t^2).
\]
This shows directly how the Hamiltonian determines
the change in the state over a small amount of time.

\paragraph{What this means physically.}
If the initial state is an energy eigenstate,
$\hat{H}\ket{E_n}=E_n\ket{E_n}$, then
\[
U(t,t_0)\ket{E_n}
=
e^{-\frac{i}{\hbar}E_n(t-t_0)}\ket{E_n}.
\]
It only gains a global phase, so its measurement
probabilities do not change.

For a superposition of energy eigenstates,
\[
\ket{\psi(t_0)}=\sum_n c_n\ket{E_n},
\]
each component gains its own phase:
\[
\ket{\psi(t)}
=
\sum_n c_n
e^{-\frac{i}{\hbar}E_n(t-t_0)}\ket{E_n}.
\]
The energy probabilities $|c_n|^2$ stay the same,
but the relative phases between components with
different energies change. These relative phases
can change interference and the probabilities of
measurements in other bases.

Finally, because $\hat{H}$ is Hermitian, the
time-evolution operator is unitary:
\[
U^\dagger(t,t_0)U(t,t_0)=\mathbb{I}.
\]
Thus, the evolution preserves normalization and
the total measurement probability remains one.

\setcounter{subsubsection}{2}
\subsubsection{The Exponential of a Hermitian Matrix}

We are given that $A$ is Hermitian, so $A^\dagger=A$.
To show that $e^{iA}$ is unitary, we need to show
\[
(e^{iA})^\dagger e^{iA}=\mathbb{I}.
\]

First, we write the exponential as a power series:
\[
e^{iA}
=
\mathbb{I}+iA+\frac{(iA)^2}{2!}
+\frac{(iA)^3}{3!}+\cdots.
\]
Taking the Hermitian conjugate conjugates the scalar
coefficients and transposes and conjugates the matrices.
Since $i^*=-i$ and $A^\dagger=A$, we have
\[
(iA)^\dagger=-iA.
\]
For powers of the same matrix,
\[
\big((iA)^n\big)^\dagger
=
\big((iA)^\dagger\big)^n
=
(-iA)^n.
\]
Although taking the Hermitian conjugate reverses the
order of a product, all the factors here are identical.

Applying this to the series gives
\[
\begin{aligned}
(e^{iA})^\dagger
&=
\mathbb{I}-iA+\frac{(-iA)^2}{2!}
+\frac{(-iA)^3}{3!}+\cdots\\
&=e^{-iA}.
\end{aligned}
\]

Now we substitute this into the unitarity condition:
\[
(e^{iA})^\dagger e^{iA}=e^{-iA}e^{iA}.
\]
The matrices $-iA$ and $iA$ commute because they are
scalar multiples of the same matrix. Therefore, we
can combine their exponentials:
\[
e^{-iA}e^{iA}
=
e^{-iA+iA}
=
e^0
=
\mathbb{I}.
\]
The reverse product gives the identity as well, so
\[
\boxed{
(e^{iA})^\dagger e^{iA}
=
e^{iA}(e^{iA})^\dagger
=
\mathbb{I}.
}
\]
Thus, $e^{iA}$ is unitary for every Hermitian matrix $A$.

Physically, this connects Hermitian Hamiltonians with
unitary time evolution. For a time-independent
Hamiltonian $H$, the operator
\[
U(t)=e^{-iHt/\hbar}
\]
has this form with $A=-Ht/\hbar$, which is Hermitian
for real $t$. Its unitarity preserves the norm:
\[
\braket{\psi(t)|\psi(t)}
=
\bra{\psi(0)}U^\dagger(t)U(t)\ket{\psi(0)}
=
\braket{\psi(0)|\psi(0)}.
\]
So a normalized initial state stays normalized, as
required for the total measurement probability to
remain one.

% Match the numbering of the exercise set.
\setcounter{section}{3}
\section{Toward Quantum Simulation}
\setcounter{subsection}{1}
\subsection{Projecting Spatial Parity}

\paragraph{Swapping the two qubits.}
Interchanging the two Hilbert spaces means swapping the
states occupying the first and second qubit positions.
We represent this with the SWAP operator:
\[
\hat{S}\ket{ab}=\ket{ba}.
\]
On the computational basis states, this gives
\[
\hat{S}\ket{00}=\ket{00},\qquad
\hat{S}\ket{01}=\ket{10},\qquad
\hat{S}\ket{10}=\ket{01},\qquad
\hat{S}\ket{11}=\ket{11}.
\]
Since the operator is linear, we can apply it separately
to each term in a Bell state.

\paragraph{The three triplet states.}
For the first state,
\[
\begin{aligned}
\hat{S}\ket{\psi_{t,1}}
&=\hat{S}\frac{\ket{00}+\ket{11}}{\sqrt{2}}\\
&=\frac{\ket{00}+\ket{11}}{\sqrt{2}}
=\ket{\psi_{t,1}}.
\end{aligned}
\]
Both terms stay the same under the swap.

For the second state,
\[
\begin{aligned}
\hat{S}\ket{\psi_{t,2}}
&=\hat{S}\frac{\ket{01}+\ket{10}}{\sqrt{2}}\\
&=\frac{\ket{10}+\ket{01}}{\sqrt{2}}
=\ket{\psi_{t,2}}.
\end{aligned}
\]
The terms exchange places, but their sum is unchanged.

For the third state,
\[
\begin{aligned}
\hat{S}\ket{\psi_{t,3}}
&=\hat{S}\frac{\ket{00}-\ket{11}}{\sqrt{2}}\\
&=\frac{\ket{00}-\ket{11}}{\sqrt{2}}
=\ket{\psi_{t,3}}.
\end{aligned}
\]
Again, each basis state stays in place. The minus sign
between the terms does not make this state antisymmetric,
because swapping the qubits does not exchange these two terms.

Therefore, all three triplet states satisfy
\[
\boxed{\hat{S}\ket{\psi_{t,j}}=+\ket{\psi_{t,j}},
\qquad j=1,2,3.}
\]
They are symmetric under exchange and have positive
spatial parity.

\paragraph{The singlet state.}
For the remaining Bell state,
\[
\begin{aligned}
\hat{S}\ket{\psi_s}
&=\hat{S}\frac{\ket{01}-\ket{10}}{\sqrt{2}}\\
&=\frac{\ket{10}-\ket{01}}{\sqrt{2}}\\
&=-\frac{\ket{01}-\ket{10}}{\sqrt{2}}
=-\ket{\psi_s}.
\end{aligned}
\]
Here, exchanging the two terms introduces an overall
minus sign. Thus,
\[
\boxed{\hat{S}\ket{\psi_s}=-\ket{\psi_s}.}
\]
The singlet is antisymmetric under exchange and has
negative spatial parity.

These results give three independent states with SWAP
eigenvalue $+1$ and one with eigenvalue $-1$, accounting
for all four dimensions of the two-qubit Hilbert space.
Here, parity refers to symmetry under exchanging the
qubits. 



\setcounter{subsection}{2}
\subsection{Time Evolution and Error Mitigation}
\setcounter{subsubsection}{0}
\subsubsection{Quantum Zeno Effect}

Our initial state is $\ket{\psi(0)}=\ket{0}$, and the
Hamiltonian is
\[
\hat{H}=\frac{X+Y}{\sqrt{2}}.
\]
Using the units in the question, the time-evolution
operator is
\[
U_{\mathrm{ext}}(t)=e^{-i\hat{H}t}.
\]
We want to find how frequently we need to measure in
the $Z$ basis to keep the state near $\ket{0}$ over
one unit of time.

\paragraph{Evolution between measurements.}
First, we calculate the square of the Hamiltonian:
\[
\hat{H}^2
=
\frac{(X+Y)(X+Y)}{2}
=
\frac{X^2+XY+YX+Y^2}{2}.
\]
Since $X^2=Y^2=\mathbb{I}$ and $XY+YX=0$, we get
\[
\hat{H}^2=\mathbb{I}.
\]
Therefore, the even powers in the exponential series
give the cosine terms, while the odd powers give the
sine terms:
\[
e^{-i\hat{H}t}
=
\mathbb{I}\cos t-i\hat{H}\sin t.
\]

After a time interval $\Delta t$, the state becomes
\[
\begin{aligned}
\ket{\psi(\Delta t)}
&=U_{\mathrm{ext}}(\Delta t)\ket{0}\\
&=\cos(\Delta t)\ket{0}
-i\frac{X+Y}{\sqrt{2}}\sin(\Delta t)\ket{0}.
\end{aligned}
\]
Using $X\ket{0}=\ket{1}$ and $Y\ket{0}=i\ket{1}$,
\[
\begin{aligned}
\ket{\psi(\Delta t)}
&=\cos(\Delta t)\ket{0}
-i\frac{1+i}{\sqrt{2}}\sin(\Delta t)\ket{1}\\
&=\cos(\Delta t)\ket{0}
+\frac{1-i}{\sqrt{2}}\sin(\Delta t)\ket{1}.
\end{aligned}
\]

\paragraph{Probability of measuring zero.}
Measuring in the $Z$ basis gives either $\ket{0}$ or
$\ket{1}$. The probability of getting $\ket{0}$ is
\[
p_0(\Delta t)
=
\left|\braket{0|\psi(\Delta t)}\right|^2
=
\cos^2(\Delta t).
\]
If the result is zero, the measurement projects the
state back to $\ket{0}$. The next interval then starts
from the same state as the previous one.

For $N$ equally spaced measurements over a total time
$T=1$, the interval is
\[
\Delta t=\frac{T}{N}=\frac{1}{N}.
\]
Here, the measurements occur at $\Delta t,2\Delta t,
\ldots,N\Delta t=1$; the initial preparation is not
counted as a measurement.

To get zero at every measurement, we multiply the
conditional probability for each interval:
\[
\boxed{
P_{\mathrm{survival}}(N)
=
\left[\cos^2\left(\frac{1}{N}\right)\right]^N.
}
\]
This counts runs in which every result is zero,
rather than runs that might leave $\ket{0}$ and
return to it before the final measurement.

\paragraph{Finding the number of measurements.}
We need
\[
\left[\cos^2\left(\frac{1}{N}\right)\right]^N
\geq 0.99.
\]
For large $N$, we can use the small-angle approximation
$\cos^2 x\approx 1-x^2$:
\[
P_{\mathrm{survival}}(N)
\approx
\left(1-\frac{1}{N^2}\right)^N
\approx
1-\frac{1}{N}.
\]
The second approximation keeps the leading correction
in the binomial expansion. This gives
\[
1-\frac{1}{N}\gtrsim 0.99,
\]
so we expect to need about $100$ measurements.

Checking the exact expression at the two nearby
integers gives
\[
P_{\mathrm{survival}}(99)
\approx 0.98994966<0.99,
\]
whereas
\[
P_{\mathrm{survival}}(100)
\approx 0.99004967>0.99.
\]
The survival probability increases with $N$ for this
equally spaced protocol. Thus, the minimum is
\[
\boxed{N_{\min}=100,\qquad \Delta t=0.01.}
\]

\paragraph{Why repeated measurements help.}
For a short interval, the probability of leaving
$\ket{0}$ is
\[
1-p_0(\Delta t)=\sin^2(\Delta t)
\approx(\Delta t)^2.
\]
Making the intervals shorter reduces this probability
quadratically. Even though we perform more measurements,
the total failure probability scales approximately as
\[
N(\Delta t)^2
=
N\left(\frac{1}{N}\right)^2
=
\frac{1}{N}.
\]
It therefore approaches zero as $N$ increases. This
is the quantum Zeno effect: frequent measurements
suppress transitions away from the initial state.

The state still evolves slightly between measurements.
For $N=100$, along a run with all zero outcomes, its
overlap probability with $\ket{0}$ stays at least
$\cos^2(0.01)\approx0.999900$ during each interval.
The measurements repeatedly return it to $\ket{0}$,
keeping it in that state's vicinity.

\setcounter{subsubsection}{1}
\subsubsection{Symmetry Identification}

Our initial state is
\[
\ket{\psi(0)}
=
\frac{\ket{00}+\ket{11}}{\sqrt{2}}
=
\ket{\Phi^+},
\]
and the time-evolution operator is
\[
U_{\mathrm{noiseless}}(t)=e^{-i(X\otimes Y)t}.
\]
We use the Bell states
\[
\ket{\Phi^\pm}
=
\frac{\ket{00}\pm\ket{11}}{\sqrt{2}},
\qquad
\ket{\Psi^\pm}
=
\frac{\ket{01}\pm\ket{10}}{\sqrt{2}}.
\]

\paragraph{(a) Evolution in the Bell basis.}
Since
\[
(X\otimes Y)^2=X^2\otimes Y^2=\mathbb{I}_4,
\]
we can rewrite the evolution operator as
\[
U_{\mathrm{noiseless}}(t)
=
\mathbb{I}_4\cos t-i(X\otimes Y)\sin t.
\]
First, we apply the Hamiltonian to the computational
basis states in $\ket{\Phi^+}$:
\[
(X\otimes Y)\ket{00}=i\ket{11},
\qquad
(X\otimes Y)\ket{11}=-i\ket{00}.
\]
Therefore,
\[
\begin{aligned}
(X\otimes Y)\ket{\Phi^+}
&=
\frac{i\ket{11}-i\ket{00}}{\sqrt{2}}\\
&=
-i\frac{\ket{00}-\ket{11}}{\sqrt{2}}\\
&=-i\ket{\Phi^-}.
\end{aligned}
\]

The state after time $t$ is then
\[
\begin{aligned}
\ket{\psi(t)}
&=U_{\mathrm{noiseless}}(t)\ket{\Phi^+}\\
&=\cos t\ket{\Phi^+}
-i(-i)\sin t\ket{\Phi^-}\\
&=\cos t\ket{\Phi^+}-\sin t\ket{\Phi^-}.
\end{aligned}
\]
Since the Bell basis is orthonormal, the measurement
probabilities are
\[
\boxed{
\begin{aligned}
P_{\Phi^+}(t)&=\cos^2 t,\\
P_{\Phi^-}(t)&=\sin^2 t,\\
P_{\Psi^+}(t)&=0,\\
P_{\Psi^-}(t)&=0.
\end{aligned}
}
\]
The probabilities sum to
\[
\cos^2 t+\sin^2 t=1.
\]
The state evolves between the two $\Phi$ Bell states,
while the two $\Psi$ Bell states remain unpopulated.

% Insert the probability plot and code here later.

\paragraph{(b) Identifying the symmetry operator.}
We need to find what the populated states have in
common. Both $\ket{\Phi^+}$ and $\ket{\Phi^-}$ contain
only $\ket{00}$ and $\ket{11}$, so the operator
$Z\otimes Z$ works.

Using $Z\ket{0}=\ket{0}$ and $Z\ket{1}=-\ket{1}$,
\[
\begin{aligned}
(Z\otimes Z)\ket{\Phi^+}
&=
\frac{\ket{00}+\ket{11}}{\sqrt{2}}
=+\ket{\Phi^+},\\
(Z\otimes Z)\ket{\Phi^-}
&=
\frac{\ket{00}-\ket{11}}{\sqrt{2}}
=+\ket{\Phi^-}.
\end{aligned}
\]
For $\ket{11}$, the two minus signs cancel.
Thus, both populated states have eigenvalue $+1$.

For the other two Bell states,
\[
\begin{aligned}
(Z\otimes Z)\ket{\Psi^+}
&=
\frac{-\ket{01}-\ket{10}}{\sqrt{2}}
=-\ket{\Psi^+},\\
(Z\otimes Z)\ket{\Psi^-}
&=
\frac{-\ket{01}+\ket{10}}{\sqrt{2}}
=-\ket{\Psi^-}.
\end{aligned}
\]
These states have eigenvalue $-1$. Therefore,
$Z\otimes Z$ distinguishes the populated states
from those that remain unpopulated:
\[
\boxed{
\langle Z\otimes Z\rangle_{\Phi^\pm}=+1,
\qquad
\langle Z\otimes Z\rangle_{\Psi^\pm}=-1.
}
\]
Here, parity refers to whether the two computational
basis bits are equal or different.

\paragraph{(c) Showing that parity is conserved.}
With $\hat{H}=X\otimes Y$, the commutator is
\[
\begin{aligned}
[Z\otimes Z,\hat{H}]
&=(Z\otimes Z)(X\otimes Y)
 -(X\otimes Y)(Z\otimes Z)\\
&=ZX\otimes ZY-XZ\otimes YZ.
\end{aligned}
\]
Since $ZX=-XZ$ and $ZY=-YZ$,
\[
\begin{aligned}
[Z\otimes Z,\hat{H}]
&=(-XZ)\otimes(-YZ)-XZ\otimes YZ\\
&=(-1)(-1)(XZ\otimes YZ)-XZ\otimes YZ\\
&=0.
\end{aligned}
\]
Therefore,
\[
\boxed{[Z\otimes Z,\hat{H}]=0.}
\]

For a time-independent observable, with $\hbar=1$,
\[
\frac{d}{dt}\langle A\rangle
=
i\langle[\hat{H},A]\rangle.
\]
Applying this to $Z\otimes Z$ gives
\[
\frac{d}{dt}\langle Z\otimes Z\rangle
=
-i\langle[Z\otimes Z,\hat{H}]\rangle
=0.
\]
The initial state has parity $+1$, so
\[
\boxed{\langle Z\otimes Z\rangle_t=+1
\quad\text{for all }t.}
\]
The state can change between $\ket{\Phi^+}$ and
$\ket{\Phi^-}$, but its parity stays the same.
This is the conserved symmetry under the time evolution.



\setcounter{subsubsection}{2}
\subsubsection{Coherent Errors: Zeno and Basis Randomization}

Now there is presence of single-qubit errors such that the time evolution becomes:
\[
U_{\mathrm{noisy}}(t)
=
e^{-iH_{\mathrm{noisy}}t},
\qquad
H_{\mathrm{noisy}}
=
X\otimes Y
+
I\otimes\frac{X+Y}{\sqrt{2}}
+
\frac{X+Y}{\sqrt{2}}\otimes I.
\]

\paragraph{(a) Plotting Evolution of Probabilities}

We use the ordered Bell basis
\[
\left\{
|\Phi^+\rangle,\,
|\Phi^-\rangle,\,
|\Psi^+\rangle,\,
|\Psi^-\rangle
\right\}.
\]
In this basis, the noisy Hamiltonian is
\[
H_{\mathrm{noisy}}
=
\begin{pmatrix}
0 & i & \sqrt{2} & 0 \\
-i & 0 & -i\sqrt{2} & 0 \\
\sqrt{2} & i\sqrt{2} & 0 & -i \\
0 & 0 & i & 0
\end{pmatrix}.
\]

Starting in $|\Phi^+\rangle$, the initial state vector is
\[
\psi(0)=
\begin{pmatrix}
1\\0\\0\\0
\end{pmatrix}.
\]
For each sampled time $t$, we numerically calculate the full
matrix exponential using Python's
\texttt{scipy.linalg.expm}, then evolve the state:
\[
\psi(t)=e^{-iH_{\mathrm{noisy}}t}\psi(0).
\]
The entries of $\psi(t)$ are the Bell-state amplitudes, so
their squared magnitudes give the measurement probabilities:
\[
P_j^{\mathrm{noisy}}(t)=|\psi_j(t)|^2.
\]

For comparison, the noiseless evolution found prior gives
\[
P_{\Phi^+}^{\mathrm{ideal}}(t)=\cos^2 t,
\qquad
P_{\Phi^-}^{\mathrm{ideal}}(t)=\sin^2 t,
\qquad
P_{\Psi^\pm}^{\mathrm{ideal}}(t)=0.
\]

Interpreting the $\pm1\%$ tolerance as an absolute probability
difference of $0.01$, we calculate
\[
\epsilon(t)
=
\max_j
\left|
P_j^{\mathrm{noisy}}(t)-P_j^{\mathrm{ideal}}(t)
\right|.
\]
All four probabilities remain within tolerance while
$\epsilon(t)\leq0.01$.

Using a time grid with spacing $\Delta t=0.0001$, the first
crossing of this tolerance occurs between $t=0.0710$ and
$t=0.0711$. Thus, the probabilities remain within tolerance
from $t=0$ until approximately
\[
\boxed{t_{\mathrm{limit}}\approx0.0710.}
\]


\part{Glossary}

% Unnumbered sections keep glossary entries separate
% from the exercise numbering.

\section*{Glossary Entries}
\addcontentsline{toc}{section}{Glossary Entries}

% Add definitions with mathematical and physical explanations.

\section*{Acknowledgments}
\addcontentsline{toc}{section}{Acknowledgments}

% Identify collaborators and their specific contributions,
% where applicable.


\end{document}